% Generated by experiments/scripts/export_latex.py (export_latex_c12.py) -- do not edit by hand.
% Source: experiments/data/2-processed-tables-and-evaluations/revised-contract-12m/discovery/discovery_tobs_both.json
% Source: experiments/data/2-processed-tables-and-evaluations/revised-contract-12m/discovery/discovery_tobs_out.json
% Source: experiments/data/2-processed-tables-and-evaluations/revised-contract-12m/discovery/discovery_t1_both.json
% Source: experiments/data/2-processed-tables-and-evaluations/revised-contract-12m/discovery/discovery_t1_out.json
\begin{table}[htbp]
\centering
\caption[계약 그래프와 발견 순서 검증의 우주의 크기.]{두 동결에서 본 계약 그래프의 크기와, 발견 순서 검증의 두 우주(크롤이 다닐 수 있는 계약 전체)의 크기예요. 합집합 그래프에는 두 층 가운데 어느 쪽에든 화살표가 있는, 서로 다른 계약의 순서 있는 쌍이 들어 있어요. 나비넥타이 분해는 합집합 그래프의 것이에요. $t_1$에서 $\mathrm{SCC}\cup\mathrm{OUT}$에는 거대 강연결 성분(SCC)의 약연결 성분(WCC)에 든 계약의 절반이 안 되는 계약만 들어 있어요. 그래서 프로토콜의 규칙에 따라, 두 동결 모두에서 우주가 WCC인 \texttt{both} 방식이 주 방식이고 \texttt{out} 방식은 민감도 분석이에요. 뜨거운 집합 $H_q(I)$에는 우주에서 0보다 큰 중요도 $I$가 가장 큰 계약들이 들어 있어요. 그 크기는 동점을 포함해 우주의 $q=1/100$이에요. 그래서 중요도가 0보다 큰 계약이 적으면 더 작아져요.}
\label{tab:discovery-sizes}
\fitwidth{%
\begin{tabular}{lrr}
\toprule
항목 & $T_{\mathrm{obs}}$ (주) & $t_1$ \\
\midrule
$G$의 계약 (두 동결에 같은 번호) & 1{,}403{,}959 & 1{,}403{,}959 \\
상한을 씌운 허락 쌍 & 115{,}006 & 17{,}940 \\
송금 쌍 & 2{,}580{,}709 & 2{,}269{,}073 \\
합집합 그래프: 쌍이 있는 계약 & 1{,}299{,}438 & 1{,}111{,}554 \\
합집합 그래프: 쌍 & 2{,}575{,}632 & 2{,}269{,}421 \\
\midrule
거대 강연결 성분 (SCC) & 265{,}804 & 222{,}592 \\
IN (SCC에 닿음) & 174{,}239 & 169{,}237 \\
OUT (SCC에서 닿음) & 242{,}833 & 109{,}804 \\
쌍이 있는 다른 계약 & 616{,}562 & 609{,}921 \\
\midrule
$|\mathrm{SCC}\cup\mathrm{OUT}|$, \texttt{out} 방식의 우주 & 508{,}637 & 332{,}396 \\
그 안의 합집합 쌍 & 1{,}546{,}906 & 1{,}283{,}946 \\
$|\mathrm{WCC}|$, \texttt{both} 방식의 우주 (주) & 1{,}169{,}044 & 1{,}015{,}867 \\
그 안의 합집합 쌍 & 2{,}378{,}418 & 2{,}121{,}895 \\
WCC 안의 상한을 씌운 허락 쌍 & 16{,}776 & 17{,}741 \\
\midrule
$|H_q(I_T)|$, \texttt{both} 방식 & 13{,}141 & 12{,}697 \\
$|H_q(I_A)|$, \texttt{both} 방식 & 11{,}662 & 8{,}136 \\
$|H_q(I_T)|$, \texttt{out} 방식 & 5{,}087 & 3{,}324 \\
$|H_q(I_A)|$, \texttt{out} 방식 & 5{,}088 & 3{,}324 \\
\bottomrule
\end{tabular}
}
\end{table}
